% @see : https://coopmaths.fr/alea/?uuid=dnb_2026_01_sujet0va_3&alea=fjpu&uuid=dnb_2024_09_metropole_4&alea=dhH9&uuid=dnb_2022_09_polynesie_4&alea=fEGA&v=eleve&pdfParam=eyJ0aXRsZSI6IiIsInJlZmVyZW5jZSI6IiIsInN1YnRpdGxlIjoiIiwic3R5bGUiOiJDb29wbWF0aHMiLCJmb250T3B0aW9uIjoiU3RhbmRhcmRGb250IiwidGFpbGxlRm9udE9wdGlvbiI6MTIsImR5c1RhaWxsZUZvbnRPcHRpb24iOjE0LCJjb3JyZWN0aW9uT3B0aW9uIjoiQXZlY0NvcnJlY3Rpb24iLCJxcmNvZGVPcHRpb24iOiJTYW5zUXJjb2RlIiwidHlwZUZpY2hlIjoiRmljaGUiLCJkdXJhdGlvbkNhbk9wdGlvbiI6IjkgbWluIiwidGl0bGVPcHRpb24iOiJTYW5zVGl0cmUiLCJuYlZlcnNpb25zIjoxLCJleG9zIjp7fX0\%3D
\documentclass[a4paper,12pt,fleqn]{memoir}
\setlrmarginsandblock{.5cm}{.5cm}{*}
\setulmarginsandblock{.5cm}{.5cm}{*}
 \usepackage{auto-pst-pdf}
%%% COULEURS %%%
\usepackage[table,svgnames]{xcolor}
\definecolor{nombres}{cmyk}{0,.8,.95,0}
\definecolor{gestion}{cmyk}{.75,1,.11,.12}
\definecolor{gestionbis}{cmyk}{.75,1,.11,.12}
\definecolor{grandeurs}{cmyk}{.02,.44,1,0}
\definecolor{geo}{cmyk}{.62,.1,0,0}
\definecolor{algo}{cmyk}{.69,.02,.36,0}
\definecolor{correction}{cmyk}{.63,.23,.93,.06}
% Couleur des filets des tableaux

\usepackage[most]{tcolorbox} % Supprimé à cause de conflits avec ProfCollege
\tcbuselibrary{documentation}

%%%%%%%% Modification paramétrage pour footnotes
% Le paquet semble déjà appelé en amont, curieux qu'il n'y ait pas eu de clash
% avec l'appel ci-dessous mais seulement quand j'ai essayé de l'inclure avec l'option lualatex

% \usepackage{hyperref}
% Parametrages
\hypersetup{
    colorlinks=true,% On active la couleur pour les liens. Couleur par défaut rouge
    linkcolor=black,% On définit la couleur pour les liens internes
    % filecolor=magenta,% On définit la couleur pour les liens vers les fichiers locaux      
    urlcolor=black,% On définit la couleur pour les liens vers des sites web
    % pdftitle={Puissance Quatre},% On définit un titre pour le document pdf
    % pdfpagemode=FullScreen,% On fixe l'affichage par défaut à plein écran
    }
% Pour avoir les numérotations arabic y compris dans les environnements tcolorbox
\renewcommand{\thempfootnote}{\arabic{mpfootnote}}
%%%%%%%%

\usepackage{multicol} 
\usepackage{calc}
\usepackage{enumitem}
\usepackage{graphicx}
\usepackage{tabularx}
%\usepackage{tablvar}

\setlength{\parindent}{0mm}
\renewcommand{\arraystretch}{1.5}
\renewcommand{\labelenumi}{\textbf{\theenumi{}.}}
\renewcommand{\labelenumii}{\textbf{\theenumii{}.}}
\setlength{\fboxsep}{3mm}
 

\setlength{\premulticols}{0pt}

\newcounter{ExoMA}

\newlength{\parindentMA}%
\setlength{\parindentMA}{\parindent} 

\usepackage{simplekv}
 

\usepackage{fancyhdr}
\fancyhead[C]{}
\fancyfoot{}

\def\bla{}
 
 


%%% MISE EN PAGE %%%

\usepackage[margin=1.5cm]{geometry}
\usepackage{setspace}
\usepackage{pgf}%
\usepackage{pgfplots}
\pgfplotsset{compat=1.18}
\usetikzlibrary{babel,arrows, arrows.meta,calc,fit,patterns,plotmarks,shapes.geometric,shapes.misc,shapes.symbols,shapes.arrows,shapes.callouts, shapes.multipart, shapes.gates.logic.US,shapes.gates.logic.IEC, er, automata,backgrounds,chains,topaths,trees,petri,mindmap,matrix, calendar, folding,fadings,through,positioning,scopes,decorations.fractals,decorations.shapes,decorations.text,decorations.pathmorphing,decorations.pathreplacing,decorations.footprints,decorations.markings,shadows}

\spaceskip=2\fontdimen2\font plus 3\fontdimen3\font minus3\fontdimen4\font\relax %Pour doubler l'espace entre les mots
 
 
  

\usepackage[french]{babel}
% loadPackagesFromContent
\usepackage{pstricks}
\usepackage{pst-plot}
\usepackage{multido}
\usepackage{pst-fun}
\newcommand{\e}{\mathrm{e}}
\usepackage{amsmath}
\usepackage{amssymb}
\usepackage{etoolbox}
\newbool{dys}
\setbool{dys}{false}          
\ifbool{dys}{
% POLICE DYS
% ===== VARIABLE =====
\def\FontChoisie{Fira} % valeurs possibles : Fira, lmodern, tgheros
\newcommand{\choiceFontsDys}[1]{
\ifstrequal{#1}{Fira}{%
  % Fira Sans + Fira Math
  % Description : Fira Sans est une police moderne, et Fira Math est une version mathématique compatible qui maintient le style sans-serif.
  % Utilisation : Fonctionne bien avec XeLaTeX et LuaLaTeX.
  \usepackage{fontspec}
  \setmainfont{Fira Sans}
  \setsansfont{Fira Sans}
  \usepackage{unicode-math}
  \setmathfont{Fira Math}
}{}
\ifstrequal{#1}{lmodern}{%
  % Latin Modern Sans
  % Description : Une version sans-serif de la célèbre police Latin Modern, qui est compatible avec mathastext.
  % Utilisation : Fonctionne avec pdfLaTeX, XeLaTeX et LuaLaTeX.
  \usepackage{lmodern}
  \renewcommand{\familydefault}{\sfdefault}
  \usepackage{mathastext}
}{}
\ifstrequal{#1}{tgheros}{%
  % TeX Gyre Heros + mathastext
  % Description : TeX Gyre Heros est une alternative à Helvetica, et le package mathastext permet d'utiliser la même police pour les mathématiques.
  % Utilisation : Fonctionne avec pdfLaTeX, XeLaTeX, ou LuaLaTeX.
  \usepackage{tgheros}
  \renewcommand{\familydefault}{\sfdefault}
  \usepackage{mathastext}
}{}
}
% Fira ou lmodern ou tgheros
\choiceFontsDys{\FontChoisie}

%\usepackage{unicode-math}
%\usepackage{fontspec}
%\setmainfont{TeX Gyre Schola}
%\setsansfont{TeX Gyre Schola}
%\setmathfont{TeX Gyre Schola Math}
%\setmainfont{OpenDyslexic}[Scale=1.0]
%\setsansfont{OpenDyslexic}[Scale=1.0]
%\setmathfont{OpenDyslexic}[Scale=1.0,range=up/{Latin,latin,num}]
\usepackage[fontsize=14]{scrextend}
\usepackage{setspace}
\setstretch{1.7}
% Une valeur d'environ 1.2 à 1.7 est souvent recommandée. Cela permet d'augmenter l'espace entre les lignes tout en offrant un léger espacement entre les mots.
\setlength{\spaceskip}{1.2em}  
% Une valeur d'environ 1.2em à 1.5em est couramment conseillée. Cela crée un espace plus ample entre les mots, ce qui peut aider à réduire la fatigue visuelle et à améliorer la fluidité de la lecture.
}{
% POLICE STANDARD
\usepackage[T1]{fontenc} 
\usepackage[scaled=1]{helvet}
\usepackage[fontsize=12]{scrextend}
}
\begin{document}
\pagestyle{empty}
% @see : https://coopmaths.fr/alea?uuid=dnb_2026_01_sujet0va_3
\subsection*{Exercice 1 : DNB Janvier 2026 Sujet 0 Version A}
On considère les fonctions $f$ et $g$ suivantes :
\[\begin{array}{l l}
f \,: &x \longmapsto 4x + 3\\
g \,:&x \longmapsto 6x
\end{array}\]
Leurs représentations graphiques $(d_1)$ et $(d_2)$ sont tracées ci-dessous :
\begin{enumerate}
\item Parmi ces deux fonctions, laquelle représente une situation de proportionnalité ?
\item Calculer l'image de 0 par la fonction $g$.
\item Déterminer l'antécédent de $0$ par la fonction $f$.

\begin{center}
\psset{xunit=1.7cm,yunit=0.325cm,arrowsize=2pt 3,comma}
\begin{pspicture*}(-1.5,-1.9)(5.5,20)
\multido{\n=-1.5+0.5}{16}{\psline[linewidth=0.25pt](\n,-1.9)(\n,20)}
\multido{\n=-0+2}{11}{\psline[linewidth=0.25pt](-1.5,\n)(5.5,\n)}
\psaxes[linewidth=1.25pt,Dx=0.5,Dy=2,labelFontSize=\scriptstyle]{->}(0,0)(-1.5,-1.9)(5.5,20)
\psplot[plotpoints=800,linewidth=1.25pt,linecolor=red]{-1.5}{4.5}{x 4 mul 3 add}
\psplot[plotpoints=800,linewidth=1.25pt,linecolor=blue]{-1.5}{4.5}{x 6 mul}
\uput[ul](3,18){\blue $(d_1)$}\uput[dr](3.75,18){\red $(d_2)$}
\end{pspicture*}
\end{center}

\item Associer à chaque droite la fonction qu'elle représente. Justifier la réponse.
\item Déterminer graphiquement les coordonnées du point d'intersection des droites $(d_1)$ et $(d_2)$.
\end{enumerate}

\bigskip

 

\subsection*{Exercice 2 : DNB Septembre 2024 Métropole}{}

On considère la fonction $f$ définie par 
\[f(x) = x^2 + 10x + 16.\]
\begin{enumerate}
\item Vérifier par le calcul que l'image de 6 par la fonction $f$ est $112$.
\item On utilise un tableur afin de calculer les images des entiers compris entre $-4$ et 4 par la
fonction $f$.

\begin{center}
\begin{tabularx}{\linewidth}{|c|*{10}{>{\centering \arraybackslash}X|}}\hline
&A&B&C&D&E&F&G&H&I&J\\ \hline
1&$x$&$-4$&$-3$&$-2$&$-1$&0&1&2&3&4\\ \hline
2&$f(x)$&$-8$&$-5$&0&7&16&27&40&55&72\\ \hline
\end{tabularx}
\end{center}

	\begin{enumerate}
		\item Parmi les 4 formules ci-dessous, recopier celle qui a été saisie dans la cellule B2, puis
étirée vers la droite afin de calculer les images des nombres donnés par la fonction $f$.

\begin{footnotesize}
\fbox{=B1*B1+10*B1+16}\, \fbox{=A1*A1+10*A1+16}\, \fbox{$=(-4)*(-4)+10*(-4)+16$} \, \fbox{$=x*x+10*x+16$}
\end{footnotesize}
		\item En utilisant le tableau, déterminer un antécédent de $0$. 
	\end{enumerate}
\item
	\begin{enumerate}
		\item Démontrer que $f(x)$ peut s'écrire $(x + 2)(x + 8)$.
		\item En déduire un autre antécédent de 0 par la fonction $f$.
	\end{enumerate}
\end{enumerate}

\bigskip



% @see : https://coopmaths.fr/alea?uuid=dnb_2022_09_polynesie_4
\subsection*{Exercice 3 : DNB Septembre 2022 Polynésie}{}


\medskip

\begin{enumerate}
\item Voici un tableau de valeurs d'une fonction $f$ :

\begin{center}
\begin{tabularx}{\linewidth}{|*{8}{>{\centering \arraybackslash}X|}}\hline
$x$		&$-2$	&$-1$	&0	&1		& 3		& 4		& 5\\ \hline
$f(x)$	&5 		&3 		&1 	&$-1$ 	&$-5$	& $-7$	& $-9$\\ \hline
\end{tabularx}
\end{center}

	\begin{enumerate}
		\item Quelle est l'image de $3$ par la fonction $f$ ?
		\item Donner un nombre qui a pour image $5$ par la fonction $f$.
		\item Donner un antécédent de $1$ par la fonction $f$.
	\end{enumerate}
\item On considère le programme de calcul suivant:
\begin{center}
\begin{tabular}{|l|}\hline
Choisir un nombre\\
Ajouter 1\\
Calculer le carré du résultat\\ \hline
\end{tabular}
\end{center}
	\begin{enumerate}
		\item Quel résultat obtient-on en choisissant $1$ comme nombre de départ? Et en choisissant $- 2$ comme nombre de départ?
		\item On note $x$ le nombre choisi au départ et on appelle $g$ la fonction qui à $x$ fait correspondre le résultat obtenu avec le programme de calcul.
		
Exprimer $g(x)$ en fonction de $x$.
	\end{enumerate}
\item La fonction $h$ est définie par $h(x) = 2x^2 - 3$.
	\begin{enumerate}
		\item Quelle est l'image de $3$ par la fonction $h$ ?
		\item Quelle est l'image de $-4$ par la fonction $h$ ?
		\item Donner un antécédent de $5$ par la fonction $h$. En existe-t-il un autre ?
	\end{enumerate}	
\item On donne les trois représentations graphiques suivantes qui correspondent chacune à une des fonctions $f$, $g$ et $h$ citées dans les questions précédentes.

Associer à chaque courbe la fonction qui lui correspond, en expliquant la réponse.

\begin{center}
\begin{tabularx}{\linewidth}{*{3}{>{\centering \arraybackslash}X}}
\psset{unit=5mm,arrowsize=2pt 3}
\begin{pspicture*}(-4,-5)(4,6)
\rput(0,5.5){Représentation \no 1}
\psaxes[labelsep=1mm,labelFontSize=\scriptstyle,linewidth=1pt,ticksize=-2pt 2pt, showorigin=false]{->}(0,0)(-4,-4.9)(4,5)
\psplot[plotpoints=1000,linewidth=1.25pt]{-2}{5}{1 2 x mul sub}
\uput[dl](0,0){\footnotesize{0}}
\end{pspicture*} &
\psset{xunit=5mm,yunit=2.5mm,arrowsize=2pt 3}
\begin{pspicture*}(-4,-4.5)(4,18.2)
\rput(0,17){Représentation \no 2}
\psaxes[labelsep=1mm,labelFontSize=\scriptstyle,linewidth=1pt,ticksize=-2pt 2pt, showorigin=false,Dy=2]{->}(0,0)(-4,-4)(4,16)
\psplot[plotpoints=1000,linewidth=1.25pt]{-3.1}{3.1}{x dup mul 2 mul 3 sub}
\uput[dl](0,0){\footnotesize{0}}
\end{pspicture*}&
\psset{xunit=4.4mm,,yunit=5mm,arrowsize=2pt 3}
\begin{pspicture*}(-5,-1.5)(3,11)
\rput(-1,10.5){Représentation \no 3}
\psaxes[labelsep=1mm,labelFontSize=\scriptstyle,linewidth=1pt,ticksize=-2pt 2pt, showorigin=false,Dx=2, subticks=1]{->}(0,0)(-5,-1.5)(3,10)
\psplot[plotpoints=1000,linewidth=1.25pt]{-4.1}{2.1}{1  x add  dup mul}
\uput[dl](0,0){\footnotesize{0}}
\end{pspicture*}\\
\end{tabularx}
\end{center}

\end{enumerate}

\bigskip



 

 
\end{document}

% Local Variables:
% TeX-engine: luatex
% End: